\subsection{残る三型と再開する具体的な入口}
\label{ns23:support-agenda}\label{ns22:support-agenda}\label{ns20:support-agenda}

再開時は最初に \nolinkurl{00_CONTINUATION_POLICY_ja.txt} を読む．
現在の高分母支持の必要候補は\emph{三型・二十六標識}である．
全候補は五頂点を使い，三点連結，最小次数三以上，八高辺以上，$b_1\ge4$ となる．
単位枝は高々二辺の matching であり，同じ core に二本が接すれば全分母一となる．
補 graph の孤立頂点を省略すると，候補は次の通りである．
\begin{center}\small
\begin{tabular}{c|c|c|c|r}
型&高辺数&次数列&低辺補 graph&標識数\\\hline
$G_{5,14}$&8&$(43333)$&$2K_2$&15\\
$G_{5,16}$&9&$(44433)$&$K_2$&10\\
$G_{5,17}$&10&$(44444)$&$\varnothing$&1
\end{tabular}
\end{center}
識別子は詰め直さない．検算は \nolinkurl{verify_support_agenda_v23.py} である．
候補数は整数接着，atlas，幾何的非同型類の存在数ではない．

\paragraph{高分母全体の原始性．}
非空 $H$ には八辺以上ある．全高分母を割る素数があるとすると，素数支持が $H$ と一致するため，
密な素数支持の定理はその素数を五，支持を $K_{2,2,1}$ に強制する．
補 graph の二辺は実際には単位分母であるが，合同分類はその和を零 modulo five にするため矛盾する．
従って高分母全体の最大公約数は一である．

二本の単位枝がある場合には任意七本の分母の最大公約数も一となり，
$G_{5,14}$ で五の倍数の分母は高々六本である．
密な $G_5=K_{2,2,1}$ と一本の実単位枝が共存する場合は $H=G_{5,16}$ に限られ，
五進支持の外のもう一本は単位枝と頂点を共有せず，分母 $q\equiv9\pmod{10}$，$q\ge9$ となる．
合同条件だけで整数 lifting を保証してはいない．

\paragraph{優先候補：二本の互いに素な低辺．}
Book と $G_{5,15}$ は全排除済みである．専用の大尺度探索や外側四高 core 完成は再開しない．
$G_{5,14}$ の二本の低辺には，三高 core 同士が接する．
第三 core の単一 pair による線形復元を両低辺に適用し，
非共鳴の分母比較と共鳴族を分けて全体の閉路条件を比較する．
固定中心四点への還元は $G_{5,15}$ の特別な構造だったため，そのまま移してはならない．
二高局所族，外側尺度表，二つの二次残余の resultant という手順のうち，
どの部分が新しい支持でも使えるかを明示する．

\paragraph{並行候補：三高同時非負性の残る格子領域．}
固定 $(s,t)$ ごとの全係数復元に，新しい四領域の排除を加える．
一般の $|\delta|\ge2$ では，尺度の正値性が許す $x,n$ の幅を先に制限する．
$\delta=0,s<t$ を標識交換だけで除けると仮定しない．canonical な $k$ の向きも変わるためである．
$\delta=1,x=0$ は $s=y=1,z\ge2$ の軸となり，固定 $(b,z)$ ごとの非零二次式で完全復元できる．
退化面 $a=0$ の大部分は固定 $s$ で有限化でき，残る $x=0,R_0=1$ の枝は
$s^2\mid2(b-1)$ を要求する．この二つの具体的な残余は別々に進められる．
反例が現れれば，全ての原方程式・奇数性・六符号・正 weights を確認して構成へ転換する．

\paragraph{より密な支持と幾何への転換．}
$G_{5,16},G_{5,17}$ では四高 core の符号 graph と元の pair・行和条件が重要になる．
三高 matching 定理を四高へ適用しない．$K_4$ を部分 graph として含むことだけでも排除できない．
整数候補が残れば全 core の weights，basin areas，integer framings，作用切片へ進む．
五次・六次 atlas の同値，不変量，別の実際の退化，open-corridor surgery への転換も，
研究全体の面白さとインパクトに応じて判断する．
